%% kapa: exam=1 round=28 year=2017 period=1 form=A subject=civil_law subject_ko=민법 seq=19 number=19
%% source: Q-Net 감정평가사 기출문제 (cst00309) artlSeq=5199603 file=extracted/1008/5199603/2017년 제28회 감정평가사 1차시험 문제지/1교시 A형.pdf
%% conversion: kordoc 4.24.51 -> tools/gen1.py (auto); review=proofread-sonnet; flags=none
\begin{question}{19}
기한의 이익에 관한 설명으로 옳은 것은? (다툼이 있으면 판례에 따름)
\begin{choices}
  \item 기한의 이익이 채권자 및 채무자 쌍방에게 있는 경우, 채무자는 기한의 이익을 포기할 수 없다.
  \item 채무자인 甲이 저당권자 乙 이외의 다른 채권자 丙에게 동일한 부동산 위에 후순위저당권을 설정해 준 경우 원칙적으로 甲은 乙에게 기한의 이익을 주장하지 못한다.
  \item 기한이익 상실의 특약은 특별한 사정이 없는 한 형성권적 기한이익 상실의 특약으로 추정된다.
  \item 형성권적 기한이익 상실의 특약이 있는 할부채무의 경우, 특별한 사정이 없는 한 1회의 불이행이 있으면 채무전액에 대하여 소멸시효가 진행한다.
  \item 정지조건부 기한이익 상실의 특약이 있는 경우, 그 특약에 정한 기한이익 상실 사유가 발생하더라도 기한이익을 상실케 하는 채권자의 의사표시가 없다면 특별한 사정이 없는 한 이행기 도래의 효과가 발생하지 않는다.
\end{choices}
\end{question}
\answer{3}
